\documentclass[11pt, a4paper]{article}

\usepackage[utf8]{inputenc}
\usepackage[T1]{fontenc}
\usepackage[ngerman]{babel}
\usepackage{amsmath}
\usepackage{amssymb}
\usepackage{graphicx}
\usepackage{geometry}
\usepackage{fancyhdr}
\usepackage{enumitem}
\usepackage{array}
\usepackage{eso-pic}

\usepackage{siunitx}
\DeclareSIUnit{\litre}{l}
\sisetup{angle-symbol-degree = ^\circ}
\sisetup{locale = DE, separate-uncertainty, inter-unit-product = {},
range-units = brackets, list-units = single, per-mode=symbol-or-fraction}

\makeatletter
\def\input@path{{./}{../../}}
\makeatother
\graphicspath{{./}{../../}}

\geometry{a4paper, top=2.5cm, bottom=2.5cm, left=2.5cm, right=2.5cm}

\input{defs}

\pagestyle{fancy}
\fancyhf{}
\cfoot{\thepage}

\begin{document}
\noindent
\vspace{9cm}
\begin{center}
    {\Huge \textbf{Lösung}} \\[1.5em]
    {\Large \textbf{Physikalische Grundlagen}} \\[1em]
    {\large 2. Schriftliche Pr\"ufung: Sommersemester 2026} \\[1em]
    {\large Studiengang: Clinical Engineering}
\end{center}

\newpage

\section*{Lösung Beispiel 1 -- Eis im Wasser}
\begin{enumerate}
    \item \textbf{Masse des warmen Wassers:}
    \[
        m_{\text{W}} = \rho_{\text{W}} V_{\text{W}}
    \]
    Erst mit Zahlenwerten:
    \[
        m_{\text{W}} = \SI{1.0}{\kilogram\per\deci\metre\cubed} \cdot \SI{3.0}{\litre} = \SI{3.0}{\kilogram}
    \]

    \item \textbf{Maximal abgebbare Wärme des warmen Wassers bis \SI{0}{\degreeCelsius}:}
    \[
        Q_{\text{ab,0}} = m_{\text{W}} c_{\text{W}} \Delta T
    \]
    \[
        Q_{\text{ab,0}} = \SI{3.0}{\kilogram} \cdot \SI{4.18}{\kilo\joule\per(\kilogram\cdot\kelvin)} \cdot \SI{40}{\kelvin} = \SI{501.6}{\kilo\joule}
    \]

    \item \textbf{Wärme zum Erwärmen und Schmelzen des Eises:}
    Zuerst wird das Eis von \SI{-10}{\degreeCelsius} auf \SI{0}{\degreeCelsius} erwärmt:
    \[
        Q_1 = m_{\text{Eis}} c_{\text{Eis}} \Delta T
    \]
    \[
        Q_1 = \SI{0.6}{\kilogram} \cdot \SI{2.09}{\kilo\joule\per(\kilogram\cdot\kelvin)} \cdot \SI{10}{\kelvin} = \SI{12.54}{\kilo\joule}
    \]
    Dann wird es geschmolzen:
    \[
        Q_2 = m_{\text{Eis}} \lambda_{\text{S}}
    \]
    \[
        Q_2 = \SI{0.6}{\kilogram} \cdot \SI{333.5}{\kilo\joule\per\kilogram} = \SI{200.1}{\kilo\joule}
    \]
    Insgesamt also
    \[
        Q_{\text{zu,0}} = Q_1 + Q_2 = \SI{12.54}{\kilo\joule} + \SI{200.1}{\kilo\joule} = \SI{212.64}{\kilo\joule}
    \]

    \item \textbf{Entscheidung und Energiebilanz:}
    Es gilt
    \[
        Q_{\text{ab,0}} = \SI{501.6}{\kilo\joule} > Q_{\text{zu,0}} = \SI{212.64}{\kilo\joule}
    \]
    Daher schmilzt das gesamte Eis. Danach wird das geschmolzene Eiswasser noch von \SI{0}{\degreeCelsius} auf die Endtemperatur $T_{\text{Misch}}$ erwärmt. Die Energiebilanz lautet somit
    \[
        m_{\text{W}} c_{\text{W}} (T_{\text{W}} - T_{\text{Misch}}) = m_{\text{Eis}} c_{\text{Eis}} (\SI{0}{\degreeCelsius} - T_{\text{Eis}}) + m_{\text{Eis}} \lambda_{\text{S}} + m_{\text{Eis}} c_{\text{W}} (T_{\text{Misch}} - \SI{0}{\degreeCelsius})
    \]

    \item \textbf{Endtemperatur:}
    Aus der Energiebilanz aus Punkt 4 folgt zunächst symbolisch
    \[
        m_{\text{W}} c_{\text{W}} (T_{\text{W}} - T_{\text{Misch}})
        = m_{\text{Eis}} c_{\text{Eis}} (- T_{\text{Eis}}) + m_{\text{Eis}} \lambda_{\text{S}} + m_{\text{Eis}} c_{\text{W}} T_{\text{Misch}}
    \]
    \[
        m_{\text{W}} c_{\text{W}} T_{\text{W}} - m_{\text{W}} c_{\text{W}} T_{\text{Misch}}
        = - m_{\text{Eis}} c_{\text{Eis}} T_{\text{Eis}} + m_{\text{Eis}} \lambda_{\text{S}} + m_{\text{Eis}} c_{\text{W}} T_{\text{Misch}}
    \]
    \[
        m_{\text{W}} c_{\text{W}} T_{\text{W}} + m_{\text{Eis}} c_{\text{Eis}} T_{\text{Eis}} - m_{\text{Eis}} \lambda_{\text{S}}
        = (m_{\text{W}} c_{\text{W}} + m_{\text{Eis}} c_{\text{W}}) T_{\text{Misch}}
    \]
    \[
        T_{\text{Misch}} = \frac{m_{\text{W}} c_{\text{W}} T_{\text{W}} + m_{\text{Eis}} c_{\text{Eis}} T_{\text{Eis}} + m_{\text{Eis}} \lambda_{\text{S}}}{(m_{\text{W}} + m_{\text{Eis}}) c_{\text{W}}}
    \]
    Erst jetzt werden die Zahlenwerte eingesetzt:
    \begin{align*}
        T_{\text{Misch}} &= \frac{\SI{3.0}{\kilogram} \cdot \SI{4.18}{\kilo\joule\per(\kilogram\cdot\kelvin)} \cdot \SI{40}{\degreeCelsius} + \SI{0.6}{\kilogram} \cdot \SI{2.09}{\kilo\joule\per(\kilogram\cdot\kelvin)} \cdot \SI{10}{\kelvin} + \SI{0.6}{\kilogram} \cdot \SI{333.5}{\kilo\joule\per\kilogram}}{(\SI{3.0}{\kilogram} + \SI{0.6}{\kilogram}) \cdot \SI{4.18}{\kilo\joule\per(\kilogram\cdot\kelvin)}} \\
        & \approx \SI{19.20}{\degreeCelsius}
    \end{align*}
\end{enumerate}

\newpage

\section*{Lösung Beispiel 2 -- Auto in der Kurve}
\begin{enumerate}
    \item \textbf{Wirkende Kräfte:}
    In vertikaler Richtung wirken Gewichtskraft $F_G$ nach unten und Normalkraft $F_N$ nach oben. In horizontaler Richtung wirkt die Haftreibungskraft zum Kurvenzentrum; sie liefert die Zentripetalkraft.

    \item \textbf{Zentripetalkraft bei \SI{54}{\kilo\metre\per\hour}:}
    Zuerst wird die Geschwindigkeit umgerechnet:
    \[
        v = \SI{54}{\kilo\metre\per\hour} = \frac{54}{3.6} \si{\metre\per\second} = \SI{15}{\metre\per\second}
    \]
    Dann gilt
    \[
        F_{\text{Zp}} = m \frac{v^2}{r}
    \]
    \[
        F_{\text{Zp}} = \SI{1200}{\kilogram} \cdot \frac{(\SI{15}{\metre\per\second})^2}{\SI{50}{\metre}} = \SI{5400}{\newton}
    \]

    \item \textbf{Kraft, die die Zentripetalkraft liefert:}
    Auf einer ebenen Stra\ss{}e liefert die Haftreibungskraft zwischen Reifen und Stra\ss{}e die notwendige Zentripetalkraft.

    \item \textbf{Maximale Haftreibungskraft:}
    Es gilt
    \[
        F_{\text{R,h,max}} = \mu_{\text{H}} F_N
    \]
    Da auf horizontaler Stra\ss{}e $F_N = mg$ ist, folgt
    \[
        F_{\text{R,h,max}} = \mu_{\text{H}} mg
    \]
    \[
        F_{\text{R,h,max}} = \num{0.55} \cdot \SI{1200}{\kilogram} \cdot \SI{9.81}{\metre\per\second\squared} = \SI{6474.6}{\newton}
    \]

    \item \textbf{Bedingung und Pr\"ufung:}
    Damit das Auto nicht rutscht, muss gelten
    \[
        F_{\text{Zp}} \leq F_{\text{R,h,max}}
    \]
    Hier ist
    \[
        \SI{5400}{\newton} \leq \SI{6474.6}{\newton}
    \]
    Die Bedingung ist also erf\"ullt. Das Auto kann die Kurve mit \SI{54}{\kilo\metre\per\hour} ohne Rutschen durchfahren.

    \item \textbf{Maximale Geschwindigkeit:}
    Im Grenzfall gilt
    \[
        F_{\text{Zp}} = F_{\text{R,h,max}}
    \]
    also
    \[
        m \frac{v_{\text{max}}^2}{r} = \mu_{\text{H}} mg
    \]
    Nach $v_{\text{max}}$ aufgel\"ost ergibt das
    \[
        v_{\text{max}} = \sqrt{\mu_{\text{H}} g r}
    \]
    Mit Zahlenwerten folgt
    \[
        v_{\text{max}} = \sqrt{\num{0.55} \cdot \SI{9.81}{\metre\per\second\squared} \cdot \SI{50}{\metre}} \approx \SI{16.42}{\metre\per\second}
    \]
    In \si{\kilo\metre\per\hour} ist das
    \[
        v_{\text{max}} \approx \SI{59.13}{\kilo\metre\per\hour}
    \]
\end{enumerate}

\newpage

\section*{Lösung Beispiel 3 -- Zylinder unter Wasser}
\begin{enumerate}
    \item \textbf{Kräfte und Kräftegleichgewicht:}
    Auf den Zylinder wirken die Gewichtskraft $F_G$ nach unten, die Auftriebskraft $F_A$ nach oben und die Federkraft $F_{\text{F}}$ ebenfalls nach oben. Im Gleichgewicht gilt daher
    \[
        F_{\text{F}} + F_A - F_G = 0
    \]

    \item \textbf{Volumen, Masse und Gewichtskraft:}
    Zuerst wird der Radius bestimmt:
    \[
        r = \frac{d}{2} = \SI{4.0}{\centi\metre} = \SI{0.04}{\metre}
    \]
    Das Volumen des Zylinders ist
    \[
        V = \pi r^2 h = \pi \cdot (\SI{0.04}{\metre})^2 \cdot \SI{0.12}{\metre} \approx \SI{6.03e-4}{\metre\cubed}
    \]
    Die Masse ergibt sich zu
    \[
        m = \rho_{\text{Al}} V = \SI{2700}{\kilogram\per\metre\cubed} \cdot \SI{6.03e-4}{\metre\cubed} \approx \SI{1.63}{\kilogram}
    \]
    Damit folgt f\"ur die Gewichtskraft
    \[
        F_G = mg \approx \SI{1.63}{\kilogram} \cdot \SI{9.81}{\metre\per\second\squared} \approx \SI{15.98}{\newton}
    \]

    \item \textbf{Auftriebskraft:}
    Nach Archimedes gilt
    \[
        F_A = \rho_{\text{W}}\cdot V \cdot g
    \]
    \[
        F_A = \SI{1000}{\kilogram\per\metre\cubed} \cdot \SI{6.03e-4}{\metre\cubed} \cdot \SI{9.81}{\metre\per\second\squared} \approx \SI{5.92}{\newton}
    \]

    \item \textbf{Formel für die Federkraft:}
    Aus der Gleichgewichtsbedingung
    \[
        F_{\text{F}} + F_A - F_G = 0
    \]
    folgt
    \[
        F_{\text{F}} = F_G - F_A
    \]

    \item \textbf{Anzeige des Federkraftmessers:}
    \[
        F_{\text{F}} = \SI{15.98}{\newton} - \SI{5.92}{\newton} \approx \SI{10.06}{\newton}
    \]
    Der Federkraftmesser zeigt also ungef\"ahr
    \[
        F_{\text{F}} \approx \SI{10.1}{\newton}
    \]
    an und wird dabei um diesen Betrag komprimiert.
\end{enumerate}

\newpage

\section*{Lösung Beispiel 4}
\begin{enumerate}
    \item \textbf{Druck}
    \item \textbf{\si{\kilogram\per\metre\cubed}}
    \item \textbf{Der Betrag der Geschwindigkeit ist konstant.}
    \item \textbf{Beim Schmelzen wird latente W\"arme ben\"otigt.}
    \item \textbf{$F = m \frac{v^2}{r}$}
    \item \textbf{Ohne resultierende Kraft bleibt ein K\"orper in Ruhe oder in gleichf\"ormig geradliniger Bewegung.}
    \item \textbf{der Druck zu.}
    \item \textbf{\SI{20}{\newton}} \quad denn
    \[
        F_A = \rho V g = \SI{1000}{\kilogram\per\metre\cubed} \cdot \SI{0.002}{\metre\cubed} \cdot \SI{9.81}{\metre\per\second\squared} \approx \SI{19.62}{\newton}
    \]
    \item \textbf{\si{\watt}}
    \item Der Wirkungsgrad
    \[
        \eta = 1 - \frac{T_{\text{kalt}}}{T_{\text{hei\ss}}}
    \]
    wird maximal, wenn $T_{\text{hei\ss}}$ möglichst groß und $T_{\text{kalt}}$ möglichst klein gewählt wird. Technische Hürden sind vor allem die begrenzte Temperaturbeständigkeit von Materialien sowie die praktische Schwierigkeit, ein sehr kaltes Reservoir dauerhaft bereitzustellen und während eines Zyklus die Te 
\end{enumerate}

\end{document}